\documentclass[12pt]{article}
\usepackage[utf8]{inputenc}
\usepackage[T1]{fontenc}
\usepackage{hyperref}
\title{Yaour\'e Language Resources for K-12 Curriculum Development}
\author{Compiled by Claude Sonnet 5.5 (claude-sonnet-5-5) for Andr\'as Kornai}
\date{}
\begin{document}
\maketitle

\section{Introduction}
Yaour\'e (ISO 639-3 yre; Glottocode yaou1238) is a Mande language of central C\^ote d'Ivoire with about 40,000 speakers (2012 figure; other sources give 24,600 or 65,000). Varieties named include Klan, Yaan, Taan, Yoo and Bhoo. It is not known to be taught in schools. The search found no primers, textbooks or dictionary for Yaour\'e, so any curriculum would have to build on the New Testament, audio recordings and the linguistic literature listed here.

\section{Primary Educational Resources}
No primer, textbook or school programme for Yaour\'e turned up in search results. The only published text in the language that was found is the New Testament (see below).

\section{Linguistic and Reference Materials}
\begin{itemize}
\item Wikipedia overview of the Yaour\'e language: \url{https://en.wikipedia.org/wiki/Yaure_language}
\item CNRS WEBBALL bibliography for Yaour\'e. It lists Hopkins's work on tonology (1982) and on the aspect-modality system (1987), Frank's phonological sketch, Egner's 1992 phonological sketches, and Gadou's 1992 study of phonological, morphological and enunciative processes. These are bibliographic entries only; I could not check whether the full texts are available: \url{https://reflex.cnrs.fr/Lexiques/webball/biblio.php?LGNOM=Yao}
\end{itemize}
No Yaour\'e dictionary was found in the search results. A Glosbe page exists, but it is a crowd-sourced translation site, not a dictionary.

\section{Religious and Community Resources}
\begin{itemize}
\item Yaour\'e New Testament (Wycliffe Bible Translators, 1999; Global Recordings Network dates it to 2005). Digital Bible Society entry: \url{https://dbs.org/bibles/YREWBT}
\item Scripture Earth listing for yre (New Testament in text, with a film of Mark and a Luke video): \url{https://www.scriptureearth.org/00eng.php?iso=yre}
\end{itemize}

\section{Digital and Audio Resources}
\begin{itemize}
\item Global Recordings Network, Yaour\'e audio recordings (evangelistic and basic Bible teaching, aimed at non-literate and oral-culture listeners): \url{https://globalrecordings.net/en/language/yre}
\end{itemize}

\section{Curriculum Development Context}
Very little educational material exists for Yaour\'e. The search found no orthography standard, no primer, no policy document and no school programme. National literacy in C\^ote d'Ivoire is a general concern, but I found no Yaour\'e-specific figure. A curriculum would need new primers and readers, drawing on the New Testament orthography and the phonological and tonal descriptions in the literature above, in cooperation with the Yaour\'e community.

\end{document}